\documentclass[11pt]{article}
\usepackage[a4paper,margin=24mm]{geometry}
\usepackage[T1]{fontenc}
\usepackage{amsmath,amssymb,amsthm,mathtools,xurl}
\usepackage[hidelinks,hypertexnames=false]{hyperref}
\allowdisplaybreaks\setlength{\emergencystretch}{3em}
\newcommand{\C}{\mathbb C}\newcommand{\R}{\mathbb R}
\newcommand{\tr}{\operatorname{tr}}\newcommand{\dd}{\dagger_G}
\newcommand{\norm}[1]{\left\lVert#1\right\rVert}
\newtheorem{theorem}{Theorem}
\title{Original correlated moments receive both native heat observations}
\author{}\date{20 September 2026 --- SZ-20260920-045}
\begin{document}\maketitle

This is a finite receiver, with a complete proof of its moment, quotient,
compression and heat maps. It continues MI1--25 \cite{MI}, NI1--19 and
QS1--32 \cite{NI}, and the finite heat action NH1--36 \cite{NH}.
It does not change result SZ-20260920-043. The companion note
\path{endpoint_metric_check/ENDPOINT_METRIC_HEAT.tex}, EM1--30,
supplies additional endpoint bounds; the bounds needed here are proved
again. The human heat-test source is Connes \cite{Connes}; only its
Gaussian choice, not its RH-assuming infinite heat theorem, provides
context. The rank-one spectral-action context is Connes--van Suijlekom
\cite{CV}. None of their infinite-dimensional results is an assumption
in this proof. All arithmetic values remain the original inputs.

\section{Original measure, relation and physical unit}

Fix the stipulated quartet, with its actual multiplicity, in the domain
\begin{equation}\tag{HR1}\label{HR1}
\begin{gathered}
0<\delta<\tfrac12,\quad\gamma>2,\quad m\ge1,\quad
k=4\ell+1\ge9,\quad c=k/2,\quad e_k=1+k(m-1),\\
g(v)=2\xi(v),\quad h(v)=\prod_{\epsilon,\eta=\pm1}
(v-\tfrac12-\epsilon\delta-i\eta\gamma)^m,\\
w_h(y)=|(g/h)(1/2+iy)|^2/(2\pi),\quad m_k=w_h^{*k},\\
\chi(S)=\prod_{a,b=0}^k
\bigl(S-c-(2a-k)\delta-i(2b-k)\gamma\bigr)^{e_k},\quad
q=e_k(k+1)^2,\quad N\ge q-1.
\end{gathered}
\end{equation}
This domain asserts no existence of an off-line zeta zero. The argument
also applies to an existing full packet closed under conjugation and
\(\rho\mapsto1-\rho\), with the original multiplicities retained,
including collisions and its complete polynomial
\(\chi=\prod_\kappa(S-\kappa)^{e_\kappa}\); the exponent is then the
maximum \(1+\sum(m_{\rho_j}-1)\) over tuples with that sum, not a count
of distinct sums. Indeed in the local tensor ring the highest nonzero
power of \(\sum z_j\), where \(z_j^{m_{\rho_j}}=0\), has degree
\(\sum(m_{\rho_j}-1)\) and nonzero multinomial coefficient on the
product of the highest powers. The next power vanishes. Taking the
least common annihilator across local summands proves this exponent.
Reflecting every tuple preserves its multiplicities and sends its sum
\(\kappa\) to \(k-\kappa=2c-\kappa\). Therefore
\(e_{2c-\kappa}=e_\kappa\) and
\(\chi(2c-S)=(-1)^q\chi(S)\). This proves the reflection quotient
map used below also in the collision case. An unpaired subpacket is
not included in this extension.

Let \(C=\C[S]/(\chi)\), in its ordered remainder basis, and let
\(M_S\) be multiplication by \(S\). Retain
\begin{equation}\tag{HR2}\label{HR2}
M=(M_S-cI)/i,\quad E_h=\C[v]/(h),\quad
\upsilon=j_h(g/h),\quad
\eta[P]=\upsilon^{\otimes k}P(v_1+\cdots+v_k).
\end{equation}
Here \(F_h\) is the original source with Mellin transform
\(\mathcal M F_h=g/h\). The constant jet of \(g/h\) at an actual zero \(\rho\) is
\(g^{(m_\rho)}(\rho)/(m_\rho!\prod_{\sigma\ne\rho}
(\rho-\sigma)^{m_\sigma})\ne0\). A finite geometric series in the
nilpotent part inverts that jet. Thus \(\upsilon\) is a unit.
The preceding annihilator calculation proves that \(\eta\) is
injective. Its marked cyclic vector is exactly
\(\eta[1]=\upsilon^{\otimes k}\), not the ambient tensor unit.
All coordinates below are carried to this image by \(\eta\), with
metric \(\langle\eta x,\eta y\rangle=x^*Gy\). Hence traces and
adjoints there are the conjugates of the displayed operators.

\section{One finite correlated arithmetic input}

Put \(P(x)=[x^2+2(\delta^2-\gamma^2)x+
(\delta^2+\gamma^2)^2]^{2m}\), of degree \(D_0=4m\), and retain
\begin{equation}\tag{HR3}\label{HR3}
e_j=\int_\R \frac{y^{2j}|g(1/2+iy)|^2}{2\pi P(y^2)}\,dy,
\quad \beta_j=\int_\R \frac{y^{2j}|g(1/2+iy)|^2}{2\pi}\,dy,
\quad x^j=Pq_j+r_j,\quad\deg r_j<D_0.
\end{equation}
Pairing the quartet factors gives \(h(1/2+iy)=
[(y-\gamma)^2+\delta^2]^m[(y+\gamma)^2+\delta^2]^m\), so these
are the original measure and moments. Polynomial division and finite
integration give
\begin{equation}\tag{HR4}\label{HR4}
e_j=\sum_{u<D_0}r_{j,u}e_u+\sum_{v\le j-D_0}q_{j,v}\beta_v,
\qquad
\zeta_j=j![z^j]\left(\sum_{u=0}^N\frac{e_u z^{2u}}{(2u)!}\right)^k
\quad(0\le j\le2N+1).
\end{equation}
The second sum is empty when appropriate; for \(j<D_0\), \(r_j=x^j\).
The coefficient at the odd index \(2N+1\) is zero exactly. The second
identity follows by expanding \((y_1+\cdots+y_k)^j\) in the original
convolution integral. All absolute moments are finite by the original
exponential envelope, so Fubini applies to each finite expansion.
In particular \(\zeta_0=e_0^k\), with no division by the mass.
The required independent scalar inputs are at most \(4m\) divided
seeds, and \(\beta_0,\ldots,\beta_{N-D_0}\) when \(N\ge D_0\).
Their certified integral evaluations and complete tails are NI1--11;
their finite division coefficients and alternate repeated-pole inputs
are QS7--30. This receiver uses their actual enclosed values, not
sampled zero parameters.

Write all primitive real inputs as one vector \(\theta\) in its
certified box \(\mathcal B\). Repeated appearances of any entry are
the same variable. Equations \eqref{HR3}--\eqref{HR4} are polynomial
in those variables and in \(\delta,\gamma\), since division is by a
monic polynomial. They therefore preserve any additional known
correlations by restricting \(\mathcal B\) to the corresponding set.
An enclosing box may include sequences not arising from measures;
it is used only for enclosure of the actual member.

For numerical radii, first set
\(\epsilon_j^e=\sum_u|r_{j,u}|\epsilon_u^{\rm seed}
+\sum_v|q_{j,v}|\epsilon_v^\beta\).
Expand the exact products to obtain the finite coefficient bound
\begin{equation}\tag{HR5}\label{HR5}
\begin{gathered}
E(z)=\sum_{u=0}^N |\widehat e_u|z^{2u}/(2u)!,\qquad
E_+(z)=\sum_{u=0}^N(|\widehat e_u|+\epsilon_u^e)z^{2u}/(2u)!,\\
|\zeta_j-\widehat\zeta_j|\le
\epsilon_j:=j![z^j](E_+^k-E^k),\qquad \epsilon_{2j+1}=0.
\end{gathered}
\end{equation}
Every product containing at least one error is present in the difference.
For enclosed coefficients, replacing a product \(av\) by its midpoint
gives error at most \(|\widehat a|\epsilon_v+
\epsilon_a(|\widehat v|+\epsilon_v)\), including the product of errors.
This proves the parameter-enclosure rule used in QS23 as well.
The midpoint sequence is obtained by evaluating the entire common map
at \(\theta_0\), not by making independent choices for repeated entries.

\section{Exact source and attained quotient, including compression}

Let \(\Pi_N\) put the even powers of \(y\) first, then the odd powers.
With indices \(0\le i,j\le N\), define
\begin{equation}\tag{HR6}\label{HR6}
H_y=[\zeta_{i+j}],\quad K_y=[\zeta_{i+j+1}],\quad
H=\Pi_NH_y\Pi_N^*,\quad K=\Pi_NK_y\Pi_N^*,\quad
J=J_{S,N}(\Pi_NT_N)^{-1},
\end{equation}
where \((T_N)_{rj}=\binom jr c^{j-r}i^r\) for \(r\le j\), and
\((T_N^{-1})_{rj}=\binom jr(-c)^{j-r}i^{-j}\).
Expansion of \((c+iy)^j\) and \(i^{-j}(S-c)^j\) proves both maps.
Here \(J_{S,N}\) is the original reduction into the physical cyclic
space, expressed in its original cyclic coordinates. In the remainder
basis its first \(q\) columns are the identity. Thus it is onto.
The parity Gram \(H\) is block diagonal; \(K\) has the two off-diagonal
parity blocks. Both receive the same \(\zeta\)'s and their literal
index shifts. Multiplication by \(y\), compressed in the source metric,
has matrix \(H^{-1}K\): its normal equations are \(Hv=Kb\).

For every positive \(H\), define
\begin{equation}\tag{HR7}\label{HR7}
G=(JH^{-1}J^*)^{-1},\quad R=H^{-1}J^*G,\quad
A=JH^{-1}KR,\quad Q=A-M,\quad T(t)=(1-t)A+tM=A-tQ.
\end{equation}
Since \(J\) is onto, \(JH^{-1}J^*>0\). Direct multiplication gives
\begin{equation}\tag{HR8}\label{HR8}
JR=I,\quad HR=J^*G,\quad R^*HR=G,\quad
GA=R^*KR=A^*G.
\end{equation}
Every lift of \(u\) is \(Ru+v\), \(Jv=0\), and its squared norm is
\(u^*Gu+v^*Hv\), since \(v^*HRu=v^*J^*Gu=0\).
This proves the full-fibre attained minimum and its unique lift.
It also proves the precise compression formula, including the first
factor \(JH^{-1}\), not a substitution of an unrelated matrix.
On the physical source retain the exact coefficient conversion
\[
P_u(S)=\sum_{j=0}^N[((\Pi_NT_N)^{-1}Ru)_j]S^j.
\]
The lift is \(P_u(\mathcal D_1+\cdots+\mathcal D_k)F_h^{\otimes k}\),
where \(\mathcal D_j=-x_j\partial_{x_j}\) and its original Mellin multiplier is the
\(j\)-th coordinate \(v_j\). Its source polynomial on the physical
line is \(P_u(c+iy)\), whose parity-ordered coefficients are exactly
\(Ru\). Moreover \(J_{S,N}(\Pi_NT_N)^{-1}Ru=JRu=u\).
Its physical image is therefore \(\eta u\), and its squared norm is
\(u^*Gu\). The parity-\(y\) coefficients have not been applied as
if they were \(S\)-monomial coefficients.

Here is the rank-one relation needed for the correlated heat tail.
Form the real monic orthogonal polynomials \(U_j(y)\), \(0\le j\le N\),
by solving their lower-degree normal equations using \(H_y>0\), and
let \(\omega_j=\langle U_j,U_j\rangle\), including
\(\omega_0=\zeta_0\). The same normal equations define the monic
\(U_{N+1}\) from the moments through \(2N+1\); no norm of that polynomial
is required. Set \(b_j=[i^jU_j((S-c)/i)]_\chi\). The coefficient map
from this orthonormal source basis to the parity source is an invertible
matrix \(B\) with \(B^*HB=I\). Hence \(BB^*=H^{-1}\),
\(E=JB\), and \(G=(EE^*)^{-1}\). Reduction of the only degree-\(N+1\)
term omitted in the source compression gives
\begin{equation}\tag{HR9}\label{HR9}
Q=r\ell,\qquad r=\frac{i}{\omega_N}b_{N+1},\quad
\ell=b_N^*G,\quad Q^2=0,\quad
\varepsilon^2:=\tr(Q^\dd Q)=\norm Q_G^2.
\end{equation}
For the sign, the full outgoing column in the orthonormal basis is
\(U_{N+1}/\sqrt{\omega_N}\). Its reduction is
\(i^{-(N+1)}b_{N+1}/\sqrt{\omega_N}\), and its lifted last-coordinate
row is \(i^N b_N^*G/\sqrt{\omega_N}\). Their product is
\(-i b_{N+1}b_N^*G/\omega_N\). Thus \(M=A-Q\), proving the sign.
Evenness gives \(U_j(-y)=(-1)^jU_j(y)\), by uniqueness of the normal
equations. The map \(\mathcal P[P(S)]=[P(2c-S)]\) satisfies
\(\mathcal PE=E\operatorname{diag}((-1)^j)\), so it is a \(G\)-isometry.
Opposite parity vectors are orthogonal; therefore \(\ell r=0\) and
\(Q^2=0\). Finally a rank-one matrix \(r\ell\) has squared operator
and Hilbert--Schmidt norm \((r^*Gr)(\ell G^{-1}\ell^*)\), by
Cauchy--Schwarz with equality in the direction \(G^{-1}\ell^*\).
This proves \eqref{HR9}, including \(Q=0\). This proof uses only the
positive finite Gram and the consistent even shifted moments, so it
holds throughout every certified positive moment box.

\section{Explicit moment radii and movement of the same fibre}

Use the consistent midpoint \(H_0,K_0\), not a diagonally shifted
matrix as if it were a moment sequence. Let
\begin{equation}\tag{HR10}\label{HR10}
e_H=\max_i\sum_{j=0}^N\epsilon_{i+j},\qquad
e_K=\max_i\sum_{j=0}^N\epsilon_{i+j+1}.
\end{equation}
The Hermitian row-sum bound gives \(\norm{H-H_0}_2\le e_H\) and
\(\norm{K-K_0}_2\le e_K\). To prove that bound, for a Hermitian
matrix \(D\), expand \(|x^*Dx|\) and use
\(2|x_i x_j|\le|x_i|^2+|x_j|^2\); the result is at most the
largest absolute row sum times \(\norm x_2^2\). Diagonalizing the
Hermitian matrix proves the operator bound.
If \(h_*>0\) is the actual source floor from MI8--9, then
\(H_0\succeq(h_*-e_H)I\). The guard \(e_H<h_*/2\) proves positivity
on the whole box and the relative bounds
\begin{equation}\tag{HR11}\label{HR11}
\rho=\frac{e_H}{h_*-e_H}<1,\quad
\sigma=\frac{e_K}{h_*-e_H},\quad
\norm{H_0^{-1/2}(H-H_0)H_0^{-1/2}}\le\rho,\quad
\norm{H_0^{-1/2}(K-K_0)H_0^{-1/2}}\le\sigma.
\end{equation}
Alternatively any directly certified midpoint floor can replace
\(h_*-e_H\). Uniform convolution-moment error \(\epsilon\) gives
\(e_H,e_K\le d\epsilon\), where
\(d=\max(\lfloor N/2\rfloor+1,\lfloor(N-1)/2\rfloor+1)\), because
only one parity block contributes in each row. Thus
\(\epsilon\le h_*/(4d)\) gives \(\rho\le1/3\).
The actual raw-input tolerance MI17 with total degree \(N\) suffices
for both \(H\) and \(K\): its apparently additional last moment
\(\zeta_{2N+1}\) is identically zero. For example
\begin{equation}\tag{HR12}\label{HR12}
\tau\le\min\{1/3,\epsilon/[3k(B+1)^{k-1}(2N)!]\},\qquad
B\ge\int\cosh y\,w_h(y)\,dy,
\end{equation}
is enough for each divided input moment. Indeed the truncated
exponential generating polynomial changes by at most \(3\tau\) on
the unit circle, its true size is at most \(B\), and telescoping
its \(k\)-th power bounds its change by \(3k\tau(B+1)^{k-1}\).
Integration against \(e^{-ij\theta}\) isolates each coefficient.
The original MI15 gives explicit \(B\), and HR5 usually improves this
common bound. The seed row sums in HR5 specify the further required
seed accuracy, with no stability assumption on their recurrence.

In the rest of this section \(J,M\) and all structural parameters
are fixed. Subscript zero means the same formulas at \(H_0,K_0\).
The attained minimum proves directly
\begin{equation}\tag{HR13}\label{HR13}
(1-\rho)G_0\preceq G\preceq(1+\rho)G_0,\qquad
\kappa=(1+\rho)/(1-\rho).
\end{equation}
Indeed bound the squared norm on every vector of the same fibre
and then take the two minima; no observation or physical unit changes.
Define, only as exact isometric proof coordinates,
\[
V=H_0^{1/2}R_0G_0^{-1/2},\quad
D=H_0^{1/2}(R-R_0)G_0^{-1/2},\quad
C_0=H_0^{-1/2}K_0H_0^{-1/2},\quad
F=H_0^{-1/2}(K-K_0)H_0^{-1/2}.
\]
Then \(V^*V=I\), \(V^*D=0\), and
\begin{equation}\tag{HR14}\label{HR14}
\norm D\le\beta:=\rho/\sqrt{1-\rho^2}.
\end{equation}
Here is a proof retaining the complete kernel of \(J\).
The columns of \(V\) are the orthogonal complement of
\(\mathcal K=\ker(JH_0^{-1/2})\), and those of \(D\) lie in \(\mathcal K\).
Put \(X=H_0^{-1/2}(H-H_0)H_0^{-1/2}\). For any target column \(y\),
write \(v=Vy\), \(w=Dy\). Then \(v^*w=0\), \(\norm v=\norm y\),
and the minimizing normal equations give
\(0=w^*(I+X)(v+w)=\norm w^2+w^*X(v+w)\).
Thus \(\norm w^2\le\rho\norm w\sqrt{\norm v^2+\norm w^2}\).
If \(w=0\) the desired bound holds; otherwise divide by \(\norm w\),
square, and rearrange to obtain
\((1-\rho^2)\norm w^2\le\rho^2\norm y^2\).
Taking the supremum proves HR14, including a zero-dimensional kernel.

Set \(Z=G_0^{-1/2}GG_0^{-1/2}\),
\(B_1=G_0^{-1/2}(R^*KR)G_0^{-1/2}\),
\(B_0=V^*C_0V=G_0^{1/2}A_0G_0^{-1/2}\), and \(a_0=\norm{A_0}_{G_0}\),
\(k_0=\norm{C_0}\). The full matrix difference is
\begin{equation}\tag{HR15}\label{HR15}
\begin{split}
B_1-B_0={}&V^*C_0D+D^*C_0V+D^*C_0D\\
&+V^*FV+V^*FD+D^*FV+D^*FD.
\end{split}
\end{equation}
In particular the two complex mixed terms involving \(C_0\) and both
complex mixed terms involving \(F\) survive in the exact identity.
Since \((V+D)^*(V+D)=I+D^*D\), grouping only for the estimate gives
\(\norm{B_1-B_0}\le k_0(2\beta+\beta^2)+\sigma(1+\beta^2)\).
Finally the coordinate matrix of \(A\) is \(Z^{-1}B_1\), and
\(\norm{Z-I}\le\rho\), \(\norm{Z^{-1}}\le(1-\rho)^{-1}\).
Consequently
\begin{equation}\tag{HR16}\label{HR16}
\norm{A-A_0}_{G_0}\le
\alpha_A:=\frac{k_0(2\beta+\beta^2)+\sigma(1+\beta^2)+\rho a_0}
{1-\rho},\quad
\norm{T(t)-T_0(t)}_{G_0}\le u:=|1-t|\alpha_A.
\end{equation}
These are finite computable bounds; they vanish as the input radii vanish.

\section{Separate trace receivers and the exact endpoint improvement}

For real \(t\), \(s>0\), retain
\begin{equation}\tag{HR17}\label{HR17}
\Phi_s(t)=\tr e^{-sT(t)^2},\quad
\Psi_s(t)=\tr e^{-sT(t)^\dd T(t)},\quad
\Delta_s(t)=\Re\Phi_s(t)-\Psi_s(t).
\end{equation}
Let \(b_0\ge\norm{T_0(t)}_{G_0}\), \(b_*=b_0+u\),
\(v=u(2b_0+u)\), and define
\(d(1)=0\), \(d(x)=(1-x^{-1})x^{-1/(x-1)}\) for \(x>1\).
Then
\begin{equation}\tag{HR18}\label{HR18}
\begin{split}
|\Phi_s(t)-\Phi_{s,0}(t)|&\le qsv\,e^{s b_*^2},\\
|\Psi_s(t)-\Psi_{s,0}(t)|&\le qsv+q d(\kappa).
\end{split}
\end{equation}
For the first line use \(\norm{T^2-T_0^2}\le v\) and the exact
Duhamel identity
\(e^{-sX}-e^{-sY}=-\int_0^s e^{-(s-u)X}(X-Y)e^{-uY}\,du\).
It follows by differentiating the product of the two exponentials;
their series justify differentiation on this compact interval.
Their norms are at most \(e^{(s-u)b_*^2}\) and \(e^{ub_0^2}\).
Finally \(|\tr F|\le q\norm F\) in orthonormal coordinates.

For the second line first compare \(T^\dagger_0T\) and
\(T_0^{\dagger_0}T_0\) in \(G_0\), where both are positive and their
difference has norm at most \(v\). The exponentials in Duhamel are
contractions, giving \(qsv\). To change the metric with \(T\) fixed,
its generalized Rayleigh quotient is
\(x^*T^*GTx/(x^*Gx)\). HR13 places it between \(\kappa^{-1}\) and
\(\kappa\) times the quotient for \(G_0\). The increasing ordered
eigenvalues satisfy the same bounds. For completeness the \(j\)-th
eigenvalue is the minimum, over \(j\)-dimensional subspaces, of the
maximum quotient there: the first \(j\) eigenvectors attain it, and
every such subspace intersects the last \(q-j+1\) eigenvectors.
This proves the comparison without identifying eigenvectors.
Both systems have \(q-r\) zero eigenvalues, \(r=\operatorname{rank}T\).
For \(x\ge0\), the maximum of \(e^{-x}-e^{-\kappa x}\) is
\(d(\kappa)\), found by differentiating and substituting
\(x=\log\kappa/(\kappa-1)\). The reverse scaled difference has
the same maximum. Thus the metric cost is at most
\(r d(\kappa)\le qd(\kappa)\), proving HR18.
Also \(d(\kappa)\le\log\kappa/e\): integrate
\(xe^u e^{-xe^u}\le1/e\) for \(0\le u\le\log\kappa\).

At \(t=1\), the sharper receiver is
\begin{equation}\tag{HR19}\label{HR19}
\begin{gathered}
\Phi_s(1)=\sum_\kappa e_\kappa e^{s(\kappa-c)^2},\qquad
\Psi_{\kappa s,0}(1)\le\Psi_s(1)\le\Psi_{s/\kappa,0}(1),\\
\Phi_s(1)-\Phi_{s,0}(1)=0,\qquad
|\Delta_s(1)-\Delta_{s,0}(1)|
=|\Psi_s(1)-\Psi_{s,0}(1)|\le \operatorname{rank}(M)d(\kappa).
\end{gathered}
\end{equation}
The symbol in the first sum denotes roots of \(\chi\); the scalar
\(\kappa\) in the bounds is the metric ratio HR13. Triangularize
the companion over \(\C\); powers and exponential series retain its
diagonal eigenvalues with all algebraic multiplicities. Alternatively
on a local block the full operator is
\(\sum_{j<e_\kappa}f_s^{(j)}(\kappa)N_\kappa^j/j!\),
\(f_s(S)=e^{s(S-c)^2}\). The nilpotent terms have zero trace, not
zero operator, and the constant term has trace
\(e_\kappa f_s(\kappa)\). This proves the first identity without
guessing or separating numerical roots. Monotonicity of the scalar
exponential and the preceding singular-value comparison prove the
squeeze. In HR1, \(M\) is invertible since \(k\) is odd and every
eigenvalue \((2b-k)\gamma-i(2a-k)\delta\) has nonzero real part.
The endpoint formula is independent of raw moments, not independent
of an uncertainty in \(\delta,\gamma\) themselves. Such structural
uncertainty is retained in the rational receiver below.

One may also use the exact isometry
\(S=(G_0^{-1}G)^{1/2}:C_G\to C_{G_0}\). It carries the physical
map and marked vector to \(\eta S^{-1}\) and \(S[1]\), so their
composition remains \(\eta[1]\). With \(D_M=SMS^{-1}-M\), the
positive square changes by
\begin{equation}\tag{HR20}\label{HR20}
E_M=M^{\dagger_0}D_M+D_M^{\dagger_0}M+D_M^{\dagger_0}D_M,
\quad |\Psi_s(1)-\Psi_{s,0}(1)|\le s\norm{E_M}_{1,G_0}.
\end{equation}
Here the trace norm is the sum of singular values in the stated
metric. To prove its use in Duhamel, write a singular-value expansion
\(F=\sum\sigma_j u_jv_j^*\); it gives
\(|\tr FB|\le\sum\sigma_j\norm B\). Both positive heat factors
are contractions. Expansion gives exactly the three terms in HR20,
including the quadratic one. In particular the endpoint radius may
be the minimum of HR19 and HR20. No change of metric scale has occurred.

\section{Finite scalar error budgets and a shared rational heat receiver}

All norms in HR11--20 may be replaced by certified upper bounds.
The following also supplies a closed sufficient raw-moment tolerance.
Let \(\epsilon\le\epsilon_*:=h_*/(4d)\). Choose constants valid for
all such centres:
\(k_0,a_0\le\bar k\), \(\norm M_{G_0}\le\bar m\).
They have the following explicit finite construction. The MI upper
source envelope gives \(H\preceq UI\); put
\(h_l=3h_*/4\), \(U_l=U+h_*/4\), \(W=(JJ^*)^{-1}\).
Then \(h_lW\preceq G_0\preceq U_lW\), so take
\(\bar m=\sqrt{U_l/h_l}\norm M_W\).
An actual absolute moment cap is
\(c_j=2A_kC_\Gamma j!/\alpha^{j+1}\) for even \(j\),
\(c_j=0\) for odd \(j\), where \(\alpha=\pi/2\).
It follows by integrating the MI2 upper exponential envelope,
discarding only the factor \((1+|y|)^{-1/2}\le1\).
Put \(V_K=\max_i\sum_j c_{i+j+1}\) and take
\(\bar k=(V_K+h_*/4)/h_l\), which bounds \(\norm{C_0}\).
Thus every constant is a finite expression in original data.

For \(\rho\le1/3\), HR16 implies
\begin{equation}\tag{HR21}\label{HR21}
\alpha_A\le C_A\epsilon,\qquad
C_A=(7\bar k+9/4)d/h_*.
\end{equation}
Indeed \(\beta\le3\rho/(2\sqrt2)\),
\(\beta^2\le3\rho/8\le1/8\), so
\(2\beta+\beta^2\le(3/\sqrt2+3/8)\rho<5\rho/2\).
Use \(a_0\le\bar k\), \((1-\rho)^{-1}\le3/2\), and
\(\rho,\sigma\le4d\epsilon/(3h_*)\).
Define \(b_c=|1-t|\bar k+|t|\bar m+|1-t|C_A\epsilon_*\).
Then HR18 and \(\log\kappa\le(9/4)\rho\le3d\epsilon/h_*\) give
\begin{equation}\tag{HR22}\label{HR22}
\begin{gathered}
R_\Phi\le C_\Phi\epsilon,\quad
C_\Phi=2qs b_c|1-t|C_A e^{s b_c^2},\\
R_\Psi\le C_\Psi\epsilon,\quad
C_\Psi=2qs b_c|1-t|C_A+3qd/(e h_*),\\
\epsilon\le\min\{\epsilon_*,\eta_\Phi/C_\Phi,
\eta_\Psi/C_\Psi\}.
\end{gathered}
\end{equation}
A zero error constant \(C_\Phi\) or \(C_\Psi\) imposes no restriction
through its corresponding tolerance quotient, and that error is zero.
The last line supplies prescribed positive moment
error budgets \(\eta_\Phi,\eta_\Psi\). For the difference use
\(\epsilon\le\eta_\Delta/(C_\Phi+C_\Psi)\). Combining HR22, HR12,
and the finite seed row sums HR5 is an explicit finite primitive-input
tolerance. It does not assign a bounded size to these constants as
\(k,N\) grow.

There is a tighter certificate that preserves dependence rather than
adding two separate input radii. For \(D\ge0\), put
\begin{equation}\tag{HR23}\label{HR23}
\begin{gathered}
P_D^\Phi(\theta)=\sum_{n=0}^D\frac{(-s)^n}{n!}\tr T(\theta)^{2n},\quad
P_D^\Psi(\theta)=\sum_{n=0}^D\frac{(-s)^n}{n!}
\tr\bigl(G(\theta)^{-1}T(\theta)^*G(\theta)T(\theta)\bigr)^n,\\
P_D^\Delta=\Re P_D^\Phi-P_D^\Psi.
\end{gathered}
\end{equation}
The same \(\theta\) is substituted everywhere before polynomial or
rational collection. Each is a rational function with real input
variables and complex rational coefficients when prescribed
\(s,t\) and box endpoints are rational and the structural parameters
are included in \(\theta\). A requested nonrational real or complex
constant can instead have its real and imaginary parts included as
additional shared real inputs with certified rational box endpoints;
the positive denominator guard must hold on that enlarged box.
Alternatively its coefficient enclosures can be refined along with
subdivision. Fixed-width independent coefficient enclosures alone are
not a claim of convergent exact-range evaluation. Complex conjugation acts on
coefficients and not on real variables. In particular all products
in HR15 and in \(T^*GT\) remain. At \(t=1\), \(P_D^\Phi\) is
identically constant in every raw-moment variable, not merely enclosed
by a small radius.

For completeness the rational denominators are explicitly positive.
Write \(h=\det H>0\), \(B_J=J\operatorname{adj}(H)J^*>0\), and
\(b=\det B_J>0\). Then
\begin{equation}\tag{HR24}\label{HR24}
G=h\,\operatorname{adj}(B_J)/b,\quad
R=\operatorname{adj}(H)J^*\operatorname{adj}(B_J)/b,\quad
G^{-1}=B_J/h.
\end{equation}
Thus products of powers of \(h,b\) give a positive common denominator
for HR23. All operations are finite determinants and products.
If \(\delta,\gamma\) are also enclosed, include those same real
variables in \(\chi,J,M\) in HR1, HR2 and HR6, rather than fixing
sample values. In the cyclic basis, the monic remainder map \(J_{S,N}\)
has polynomial coefficients and is onto throughout, by its identity
minor. With \(\delta_->0,\delta_+<1/2,\gamma_->2\) and a verified
source floor over the box, HR24 proves positivity throughout.
The scalar movement formulas HR13--22 are stated for fixed \(J,M\);
the rational formulas HR23--24 require no such freeze and enclose the
whole shared structural family. The physical identification is asserted
at the original actual input, not at nonzero test parameters masquerading
as zeros.

Here is an exact box-enclosure algorithm, including its termination
for every required positive denominator guard. Affinely map each
finite box to \([0,1]^r\). For \(P(x)=\sum_\beta a_\beta x^\beta\)
with coordinate degrees \(d_i\), its tensor Bernstein coefficients are
\begin{equation}\tag{HR25}\label{HR25}
b_\alpha=\sum_{\beta\le\alpha}a_\beta
\prod_i\binom{\alpha_i}{\beta_i}/\binom{d_i}{\beta_i}.
\end{equation}
To prove the formula, use the binomial identity
\(x^j=\sum_{l=j}^d\binom lj\binom dj^{-1}\binom dl
x^l(1-x)^{d-l}\) in each variable; expanding \((x+1-x)^{d-j}\)
proves that identity. The nonnegative basis sums to one, so coefficient
extrema enclose the polynomial. Apply this to numerator and positive
denominator, and divide the resulting intervals, using all four endpoint
quotients. For a complex numerator do this for real and imaginary parts.
Subdivide rationally if a denominator lower bound is not yet positive.
For the exact fixed polynomial coefficients just specified, this
terminates on a compact positive box: on a subbox of side at most
\(\delta\), expansion about its corner has constant coefficient the
value there and every other monomial coefficient bounded by
\(C\delta\), with a finite \(C\) from the original polynomial
coefficients. Formula HR25 bounds each Bernstein coefficient's
deviation from that value by another explicit finite constant times
\(\delta\). The positive continuous denominator has positive minimum
on the compact box. Thus all sufficiently small subboxes certify it.
The same estimate proves convergence of the enclosure widths to the
actual range as the mesh tends to zero. This is a finite algorithm
for any prescribed excess over that range, not a promise that a genuine
input uncertainty interval collapses to a point.

\section{Complete heat tails, with correlation and outgoing curvature}

Choose uniform bounds \(b\ge\norm{T(\theta)}_{G(\theta)}\) and
\(E_Q\ge\norm{Q(\theta)}_{G(\theta)}^2\) on the box.
For fixed structural data explicit choices are
\begin{equation}\tag{HR26}\label{HR26}
b=\sqrt\kappa(b_0+u),\qquad
E_Q=\kappa(\norm{A_0-M}_{G_0}+\alpha_A)^2.
\end{equation}
For variable structural data they can instead be obtained from the
positive rational forms, for example by bounding \(G,G^{-1},T,Q\)
entrywise on the compact box. The inequality
\(\norm X_G\le\sqrt{\norm G\norm{G^{-1}}}\norm X_2\), together
with \(\norm X_2\le\norm X_F\), proves these finite upper bounds.
Put \(x=sb^2\). The separate rigorous truncation errors are
\begin{equation}\tag{HR27}\label{HR27}
|\Phi_s-P_D^\Phi|\le q e^x\frac{x^{D+1}}{(D+1)!},\qquad
|\Psi_s-P_D^\Psi|\le q\frac{x^{D+1}}{(D+1)!}.
\end{equation}
For the first, the matrix exponential tail is bounded term by term;
\(\sum_{j>D}x^j/j!\le e^x x^{D+1}/(D+1)!\), because
\((D+1+l)!\ge(D+1)!l!\). For the second, diagonalize the positive
\(T^\dd T\). For a scalar \(y\in[0,x]\), Taylor's integral remainder
for \(e^{-y}\) has absolute value at most \(y^{D+1}/(D+1)!\), since
all derivative absolute values on \([0,y]\) are at most one. Summing
over all \(q\) eigenvalues proves the second bound, with zeros included.

The correlated tail is substantially different. Define
\(D_n=\tr(T^\dd T)^n-\Re\tr T^{2n}\). Then, throughout the box,
\begin{equation}\tag{HR28}\label{HR28}
D_1=t^2\varepsilon^2,\qquad
|D_n|\le n(2n-1)t^2\varepsilon^2 b^{2n-2},\qquad
\Delta_s=\sum_{n\ge1}\frac{(-1)^{n+1}s^n}{n!}D_n.
\end{equation}
Here is the complete finite proof of the dimension-free bound. Write
\(T=U+iV\), \(U=(T+T^\dd)/2\), \(V=(T-T^\dd)/(2i)\), both
selfadjoint. HR9 gives \(\tr V^2=t^2\varepsilon^2/2\).
For \(Z(v)=U+ivV=((1+v)T+(1-v)T^\dd)/2\), \(0\le v\le1\),
both \(\norm{Z(v)}_G\) and its adjoint norm are at most \(b\).
Set \(f_n(v)=\tr(Z(v)^\dd Z(v))^n-\Re\tr Z(v)^{2n}\).
At zero it vanishes; its first derivative vanishes because the alternating
word has equal numbers of \(+iV\) and \(-iV\), whose traces cancel by
cyclicity, while the other derivative is the real part of
\(2ni\tr(U^{2n-1}V)\), zero since the trace in it is real.
A second derivative of either word has \(2n(2n-1)\) ordered terms,
each with two specified factors \(V\). A cyclic word \(VX VY\) has
absolute trace at most \(\norm V_{\rm HS,G}^2\norm X_G\norm Y_G\).
This follows from matrix-entry Cauchy--Schwarz in the exact coordinates
\(G^{1/2}\), and from summing squared column norms for multiplication
of a Hilbert--Schmidt matrix by a bounded matrix. Hence
\(|f_n''|\le2(2n)(2n-1)\tr(V^2)b^{2n-2}\).
Integrate with weight \(1-v\), whose integral is \(1/2\), and use
\(\tr V^2=t^2\varepsilon^2/2\). This proves the bound in HR28.
Direct expansion gives \(D_1=2\tr V^2\). Absolute convergence of the
two exponential series proves the last identity.

For every integer \(D\ge1\), summing that bound yields the explicit
correlated certificate
\begin{equation}\tag{HR29}\label{HR29}
\boxed{\quad |\Delta_s-P_D^\Delta|
\le s t^2 E_Q e^x(2D+1)\frac{x^D}{D!}.\quad}
\end{equation}
Indeed, with \(j=n-1\), the remaining scalar sum is
\(s t^2E_Q\sum_{j\ge D}(2j+1)x^j/j!\), equal to
\(s t^2E_Q[2x\sum_{j\ge D-1}x^j/j!+\sum_{j\ge D}x^j/j!]\).
Apply the preceding factorial tail bound to both sums. The formula
also holds at \(b=0\) and \(t=0\), with zero error and no division
by either. For \(D=0\), the bound is
\(s t^2E_Q(1+2x)e^x\). All bounds tend to zero as \(D\to\infty\)
for the fixed finite input box: the ratio of successive factorial
terms is eventually at most any chosen constant below one.

Thus a rational interval \([L_D,U_D]\) for the single collected
\(P_D^\Delta\) from HR23--25 gives the certified final interval
\begin{equation}\tag{HR30}\label{HR30}
\Delta_s(\theta)\in[L_D-\mathcal R_D,U_D+\mathcal R_D],\qquad
\mathcal R_D=s t^2E_Qe^x(2D+1)x^D/D!.
\end{equation}
Add any outward coefficient-arithmetic error once, not independently
for each repeated occurrence. Separate trace intervals use HR27.
At the endpoint one may use the exact full-root HR19 for \(\Phi_s\)
and the positive trace polynomial only; no polynomial roots need be
numerically located. The common difference polynomial preserves its
exact factor \(t^2\), since each \(D_n\) is polynomial in \(t\) with
zero constant and linear coefficients by the proof above.

Finally the same estimate at degree one gives
\begin{equation}\tag{HR31}\label{HR31}
|\Delta_s-st^2\varepsilon^2|
\le st^2\varepsilon^2[(1+2x)e^x-1],\quad
\varepsilon^2=\lim_{s\downarrow0}\Delta_s/(st^2)\quad(t\ne0).
\end{equation}
The first bound uses the exact scalar tail from \(j=1\), and the
limit follows since its bracket tends to zero. For \(0\le x<1\),
the coefficientwise bound \(e^x\le(1-x)^{-1}\) makes that bracket
at most \(3x/(1-x)\). In particular \(x\le1/7\) gives relative
error at most \(1/2\). The result supplies a finite receiver of
actual arithmetic inputs for both observations and their difference;
it asserts no unevaluated centres, signs, or degree-uniform decay.
Letting \(N,k\) grow changes all displayed constants and is not a
limit licensed by any estimate here.

\section{Executed exact comparison, from one shared input to the final interval}

The exact checker uses comparison measures, explicitly not arithmetic
zero packets. Its end-to-end family is
\(\nu_w=\delta_{-1}+w\delta_0+\delta_1\), \(1\le w\le2\),
with third convolution, original sum coordinate, \(N=2\), and comparison
relation \(y^2+1\). This is outside HR1's arithmetic domain and tests
the finite receiver, not the existence of a quartet. Put \(m_0=w+2\).
Its moments are
\(\zeta_0=m_0^3\), \(\zeta_2=6m_0^2\),
\(\zeta_4=6m_0^2+72m_0\), and all odd moments vanish.
These follow by taking coefficients in
\((m_0+z^2+z^4/12)^3\) with their factorials.
In parity order \((1,y^2,y)\), the quotient map is
\(J=\left(\begin{smallmatrix}1&-1&0\\0&0&1\end{smallmatrix}\right)\).
Direct substitution in HR7 gives
\begin{equation}\tag{HR32}\label{HR32}
G=\operatorname{diag}\!\left(\frac{6m_0^3}{m_0+12},6m_0^2\right),
\quad A=\begin{pmatrix}0&-1\\-m_0/(m_0+12)&0\end{pmatrix},
\quad M=\begin{pmatrix}0&-1\\1&0\end{pmatrix}.
\end{equation}
For example the determinant of the even source block is
\(6m_0^4(m_0+6)>0\); its quadratic inverse on \((1,-1)\) is
\((m_0+12)/(6m_0^3)\), proving the first quotient entry. The odd
entry is \(\zeta_2\), and multiplication in HR7 gives the stated \(A\).
Writing \(r=(m_0+12)/m_0\), the endpoint satisfies
\(\Phi_s=2e^s\), \(\Psi_s=e^{-sr}+e^{-s/r}\), and
\(\varepsilon^2=r+r^{-1}+2\). Since \(3\le m_0\le4\),
\(4\le r\le5\). In particular \(b^2\le10\) and
\(\varepsilon^2\le16\), bounds deliberately rounded upwards and
also checked by exact Bernstein inequalities in the script.

The executed certificate uses \(s=1/100,D=3\), and ten rational
subintervals of \([1,2]\). It collects the shared difference HR23
before enclosing it. Its polynomial interval is
\[
\left[\frac{7451308721303667}{131072000000000000},
\frac{10722634835189641}{136687500000000000}\right].
\]
HR29 and \(e^{1/10}\le10/9\) give the rational tail \(7/33750\).
The resulting complete heat-difference interval is
\[
\left[\frac{200451332275199009}{3538944000000000000},
\frac{10750984835189641}{136687500000000000}\right].
\]
Every numerator, denominator and radius is obtained by rational
arithmetic. Additional tests retain nonunit source mass, the original
complex \(S\)-coordinate coefficient maps, a nontrivial quotient kernel,
and the defective relation \((y^2+1)^2\), including both multiplicities.
The normal and optimized runs use explicit exception checks, not removable
assertions. Their receipts identify the proof and checker by SHA-256.

\begin{thebibliography}{9}
\bibitem{MI} Split-Zero programme, \emph{Moving-interval native margins
and finite raw-moment precision}, 20 September 2026,
\texttt{MOVING\_INTERVAL\_NATIVE\_PRECISION.tex}, MI1--25,
SZ-20260920-038. Complete proof read for this receiver.
\bibitem{NI} Split-Zero programme, \emph{Evaluating the original
quartet-divided moments and carrying their errors through both minima},
\texttt{NATIVE\_INPUT\_EVALUATION.tex}, NI1--19, and
\emph{Exact finite input reduction for the original quartet-divided moments},
\texttt{QUARTET\_MOMENT\_SEED\_REDUCTION.tex}, QS1--32,
20 September 2026, SZ-20260920-041. Both complete proofs read.
NI's original input integration uses F. Johansson,
\emph{Numerical integration in arbitrary-precision ball arithmetic},
arXiv:1802.07942v1, sections 2--3; that author source was read by the
NI owner, not independently here. This note proves its receiver and
does not claim new executions of that integral provider.
\bibitem{NH} Split-Zero programme, \emph{The finite native heat action},
\texttt{NATIVE\_HEAT\_ACTION.tex}, NH1--36, and
\texttt{metric\_check/HEAT\_METRIC\_COMPARISON.tex}, HM1--19,
20 September 2026, SZ-20260920-043. Complete proofs read and retained
unchanged. Their receiver identities are proved here in raw coordinates.
\bibitem{Connes} Alain Connes, \emph{Heat expansion and zeta},
\url{https://arxiv.org/abs/2402.13082v1}, original author source
\texttt{heat-kernel-zeta.tex}, Gaussian formula \texttt{ft}.
Complete original source read by the preceding heat worker, as recorded
in NH; not independently reread in this continuation. Its theorem
\texttt{heatkernelintro} assumes RH and is not used.
\bibitem{CV} Alain Connes and Walter D. van Suijlekom,
\emph{Quadratic Forms, Real Zeros and Echoes of the Spectral Action},
\url{https://arxiv.org/abs/2511.23257v1}, original author source
\texttt{Araki-final-oct25.tex}, section \texttt{sect:sa}.
Context and source reading boundary inherited from NH; no new source
inspection is claimed. All finite estimates used here are proved above.
\end{thebibliography}
\end{document}
